% !TEX root = Bachelorarbeit Synthetische Daten.tex
\begin{titlepage}
  % Andere Seitenränder als im Rest der Arbeit
  \newgeometry{lmargin=2cm,tmargin=10mm,rmargin=15mm,bmargin=1cm}
  % Hausfarbe der HAW; diese und die folgenden Einstellungen sind lokal und gelten nur innerhalb der Umgebung "titlepage"
  \color{haw}
  % Blocksatz für die Titelseite deaktivieren
  \raggedright
  % Logos rechtsbündig setzen
  \hspace{\fill}\includegraphics[width=5.00cm]{logo_haw.jpg}\\
  \vspace{0.4cm}
  \hspace{\fill}\includegraphics[width=4.30cm]{logo_fraunhofer-ipk.png}\\

  \vspace{5cm}

  % Hausschrift Open Sans der HAW
  \setmainfont{Open Sans}
  % Etwas kleiner als üblich
  \small
  % Fett und in Majuskeln
  \textbf{BACHELORARBEIT}

  \vspace{8mm}

  % Titel der Arbeit
  \begin{minipage}{0.8\linewidth}
    % Haupttitel
    \setmainfont{Martel Heavy} \LARGE
    Contrastive Learning\\[1mm] % [1mm] für etwas größeren Durchschuss
    mit Stable Diffusion-basierter\\[1mm]
    Datenaugmentation\\[4mm]
    % Untertitel
    \setmainfont{Open Sans} \Large
    Verbesserung der Bildklassifikation\\[1mm]
    durch synthetische Daten\\[1mm]
    % Horizontaler Strich
    \,\rule{11mm}{1.2mm}
  \end{minipage}

  \vspace{9.2mm}

  vorgelegt am 16. September 2024\\
  Paul Hofmann

  \vspace{3cm}

  % Prüfer*innen & Anschriften
  \hspace*{37mm}
  \begin{minipage}{0.5\linewidth}
    \begin{tabular}{@{}ll}
      Erstprüferin: & Prof. Dr. Larissa Putzar\\[-.3mm]
      Zweitprüfer: & Prof. Dr. Jan Neuhöfer \\
    \end{tabular}\\
	
	% Horizontaler Strich
    %\,\rule{9mm}{1mm}\\[1.5mm]
    \vspace{5mm}
	
	% Anschrift HAW
	\begin{tabular}{@{} p{0.75\linewidth} ll} %{p{0.3\linewidth}p{0.3\linewidth}}
    \textbf{HOCHSCHULE FÜR ANGEWANDTE} & \textbf{FRAUNHOFER-INSTITUT FÜR}\\
    \textbf{WISSENSCHAFTEN HAMBURG} & \textbf{PRODUKTIONSANLAGEN UND}\\
    Department Medientechnik & \textbf{KONSTRUKTIONSTECHNIK IPK}\\
    Finkenau 35 & Pascalstraße 8–9\\
    22081 Hamburg & 10587 Berlin
	\end{tabular}\\
  \end{minipage}
\end{titlepage}
% Setzt Geometrie wieder auf Standardwerte zurück
\restoregeometry

% Seite mit dem Abstract ohne Seitenzahl ausgeben
\thispagestyle{empty}

% Abstract
\section*{Zusammenfassung}

Diese Bachelorarbeit untersucht die Verbesserung der Bildklassifikation durch synthetische Daten im Supervised Contrastive Learning (SCL). Ziel ist es, die Eignung von DA-Fusion, einer Stable Diffusion-basierten Methode zur Datenaugmentation, für die Generierung synthetischer Daten in einem Anwendungsfall aus der Recyclingwirtschaft zu evaluieren. Darüber hinaus wird untersucht, wie generierte Near Out-of-Distribution (OOD)-Daten in das Supervised Contrastive Learning integriert werden können, um als negativ-Beispiele zu dienen und die Repräsentationen der In-Distribution-Daten weiter zu verbessern. Die Ergebnisse zeigen, dass DA-Fusion geeignete In-Distribution-Augmentationen liefert, welche die Klassifikationsleistung verbessern. Jedoch führte die Integration von Near OOD-Daten im SCL zu einer Verschlechterung der Modellgenauigkeit, da die Beispiele oft zu weit von den In-Distribution-Daten entfernt waren. Die Arbeit hebt die Herausforderungen bei der Generierung und Nutzung synthetischer Daten hervor und bietet Ansätze zur Optimierung von Contrastive Learning in diesem Kontext.

% In Englisch
{
  \begin{otherlanguage}{english}
    \section*{Abstract}

      This bachelor thesis investigates how image classification can be improved using synthetic data for Supervised Contrastive Learning (SCL). The aim is to evaluate the suitability of DA-Fusion, a stable diffusion-based data augmentation method, for generating synthetic data in a use case from the recycling industry. It also investigates how synthetic near out-of-distribution (OOD) images can be integrated into supervised contrastive learning to serve as negative examples in order to further improve the representation of in-distribution data. The results show that DA-Fusion provides suitable in-distribution augmentations that improve classification performance. However, the integration of near OOD data in SCL led to a deterioration of model accuracy, as the examples were often too dissimilar from the in-distribution data. The thesis highlights the challenges of generating and using synthetic data and offers approaches to optimise contrastive learning in this context.
  \end{otherlanguage}
}
